\documentclass[../../../include/open-logic-section]{subfiles}

\begin{document}

\olfileid{sth}{cardinals}{zfc}	
\olsection{$\ZFC$: ఒక మైలురాయి}

సుక్రమపరచడం స్వయంసిద్ధాన్ని చేర్చడంతో చివరి సిద్ధాంతపరమైన మైలురాయిని చేరుకున్నాం. ఇప్పుడు~$\ZFC$కు కావలసిన స్వయంసిద్ధాలన్నీ మన దగ్గర ఉన్నాయి. వివరంగా:

\begin{defn}
$\ZFC$ సిద్ధాంతంలో ఇవి స్వయంసిద్ధాలు: విస్తరణత, సమ్మేళనం, యుగ్మాలు, ఘాత సమితులు, అనంతత్వం, పునాది, సుక్రమపరచడం; అలాగే విభజన, ప్రతిస్థాపన పథకాల అన్ని ఉదాహరణలు. మరొక విధంగా చెప్పాలంటే,~$\ZF$కు సుక్రమపరచడాన్ని చేర్చినదే $\ZFC$.
\end{defn}

$\ZFC$ అంటే \emph{ఎంపిక}తో కూడిన \emph{జెర్మెలో--ఫ్రెంకెల్} సమితి సిద్ధాంతం. మనం చేర్చిన స్వయంసిద్ధం పేరు ``ఎంపిక'' కాక ``సుక్రమపరచడం'' కాబట్టి ఇది కొంచెం విచిత్రంగా అనిపించవచ్చు. కానీ తరువాత \emph{ఎంపిక}ను సూత్రీకరించినప్పుడు, సుక్రమపరచడం $\ZF$ పరిధిలో ఎంపికకు సమానమని తేలుతుంది (\olref[choice][woproblem]{thmwochoice} చూడండి). కాబట్టి రెండింటిలో దేన్ని ``ప్రాథమిక'' స్వయంసిద్ధంగా తీసుకున్నా తేడా లేదు. సాహిత్యంలో ``$\ZFC$'' అనే పేరే ప్రామాణికం.

\end{document}
